%% ============================================================
%%  TRABAJO AUTÓNOMO N.º 3
%%  Maestría en Gestión y Administración Financiera Pública
%%  Formato académico en economía
%% ============================================================

\documentclass[12pt, a4paper]{article}

% --- Codificación y lenguaje ---
\usepackage[utf8]{inputenc}
\usepackage[T1]{fontenc}
\usepackage[spanish,es-tabla]{babel}
\usepackage{lmodern}

% --- Márgenes ---
\usepackage[top=2.5cm, bottom=2.5cm, left=3cm, right=3cm]{geometry}

% --- Matemáticas ---
\usepackage{amsmath, amssymb, amsthm}

% --- Tablas ---
\usepackage{booktabs}
\usepackage{multirow}
\usepackage{longtable}
\usepackage{array}

% --- Figuras ---
\usepackage{graphicx}
\usepackage{subcaption}
\graphicspath{{figuras/}}

% --- Colores e hipervínculos ---
\usepackage[colorlinks=true,
            linkcolor=black,
            citecolor=blue,
            urlcolor=blue]{hyperref}

% --- Listas ---
\usepackage{enumitem}

% --- Notas al pie ---
\usepackage{footmisc}

% --- Formato de captions ---
\usepackage{caption}

% ============================================================
%  MACROS
% ============================================================

% Formato para montos en guaraníes.
% Se utiliza \text{} para evitar conflictos con babel
% cuando el macro aparece dentro de modo matemático.
\newcommand{\Gs}[1]{Gs.\ \text{#1}}

% ============================================================
%  ENTORNOS MATEMÁTICOS
% ============================================================

\theoremstyle{plain}
\newtheorem{theorem}{Teorema}[section]
\newtheorem{proposition}[theorem]{Proposición}
\newtheorem{lemma}[theorem]{Lema}
\newtheorem{corollary}[theorem]{Corolario}

\theoremstyle{definition}
\newtheorem{definition}{Definición}[section]
\newtheorem{assumption}{Supuesto}

\theoremstyle{remark}
\newtheorem{remark}{Observación}

% ============================================================
%  METADATOS DEL DOCUMENTO
% ============================================================

\title{%

    \textbf{Maestría en Gestión y Administración Financiera Pública}\\[6pt]
    {\large Auditoria Gubernamental MECIP}\\\

    \large Trabajo Autónomo\\\
    
    \large Modulo: 
    
}

\author{%
    Aqui va el Autor\\[6pt]
}

\date{2026}

% ============================================================
\begin{document}
% ============================================================

\maketitle
\thispagestyle{empty}

\newpage
\tableofcontents
\newpage

% ============================================================
\section{Introducción y objetivo}
% ============================================================

De acuerdo con la consigna del trabajo práctico, se toma como base el listado
de comprobantes suministrado (\texttt{1393.csv}) y se elabora el muestreo
aplicando la fórmula que corresponde según el tipo de variable a auditar,
explicando al final la ficha técnica resultante.

El listado corresponde a un libro de ventas: cada fila es un comprobante
(factura de contado, factura de crédito o nota de crédito) con su fecha,
número, timbrado, RUC/CI del cliente, montos gravados, IVA y total. El
objetivo de auditoría es examinar una muestra representativa de esos
comprobantes para validar la existencia y exactitud de los importes
registrados.

% ============================================================
\section{Marco conceptual}
% ============================================================

\subsection{Población y muestra}

La \textbf{población} es el conjunto finito o infinito de elementos con
características comunes sobre el cual se harán extensivas las conclusiones
de la revisión; la \textbf{muestra} es el subconjunto representativo que se
extrae de la población accesible \cite{cofae}. El proceso para seleccionar
esa muestra es el \textbf{muestreo}.

\subsection{Modelos estadísticos de tamaño de muestra}

El material presenta dos modelos, según el tipo de variable que se quiere
medir:

\begin{enumerate}
    \item \textbf{Fórmulas según Kish} --- para estimar una
    \emph{proporción o atributo} (por ejemplo, qué porcentaje de comprobantes
    presenta un determinado incumplimiento).

    \item \textbf{Estimación de la medida por unidad} --- para estimar o
    poner a prueba un \emph{monto} (una variable continua), como el importe
    total de los comprobantes.
\end{enumerate}

\subsubsection{Fórmulas según Kish}

\begin{align}
    n' &= \frac{s^2}{V^2} \\
    n &= \frac{n'}{1 + n'/N} \\
    V^2 &= S_e^2 \\
    s^2 &= \rho (1-\rho)
\end{align}

donde $N$ es el tamaño de la población, $\rho$ la probabilidad de ocurrencia
esperada del atributo, $S_e$ el error estándar deseado, $n'$ la muestra sin
ajustar y $n$ la muestra ajustada por población finita.

\subsubsection{Estimación de la medida por unidad}

Cuando la variable de interés es un monto, el tamaño de muestra se calcula
a partir de la desviación estándar del monto y del error tolerable definido
para la auditoría:

\begin{align}
    TPRM &= \frac{ET}{1 + \dfrac{CAI}{CRI}} \\[4pt]
    n &= \frac{N \times CRI \times De}{TPRM}
\end{align}

donde:
\begin{itemize}[nosep]
    \item $N$ = tamaño de la población a muestrear;
    \item $ET$ = error tolerable (el error monetario máximo que puede existir
          en la partida sin afectar de forma material la información que
          la integra);
    \item $CAI$, $CRI$ = coeficientes de aceptación y de rechazo incorrecto,
          según el nivel de riesgo aceptado por el auditor (Tabla
          \ref{tab:riesgo});
    \item $De$ = desviación estándar estimada del monto;
    \item $TPRM$ = tolerancia planeada para el riesgo de muestreo.
\end{itemize}

\begin{table}[htbp]
    \centering
    \caption{Coeficientes de aceptación y rechazo incorrecto según el nivel
    de riesgo aceptable}
    \label{tab:riesgo}
    \begin{tabular}{@{}ccc@{}}
        \toprule
        Nivel de riesgo aceptable (\%) & CAI & CRI \\
        \midrule
        1{,}00  & 2{,}33 & 2{,}58 \\
        4{,}60  & 1{,}68 & 2{,}00 \\
        \textbf{5{,}00}  & \textbf{1{,}64} & \textbf{1{,}96} \\
        10{,}00 & 1{,}28 & 1{,}44 \\
        15{,}00 & 1{,}04 & 1{,}28 \\
        20{,}00 & 0{,}84 & 1{,}15 \\
        25{,}00 & 0{,}67 & 1{,}04 \\
        \bottomrule
    \end{tabular}
\end{table}

\subsection{Distribución de la muestra por estrato}

Cuando la población está dividida en estratos (por ejemplo, tipos de
documento), la muestra se reparte proporcionalmente mediante un factor
multiplicador:

\begin{equation}
    f = \frac{n}{N}, \qquad n_h = \lceil N_h \times f \rceil
\end{equation}

donde $N_h$ es el tamaño de cada estrato y $n_h$ la cantidad de elementos a
seleccionar dentro de ese estrato.

\subsection{Modelo de riesgo de auditoría}

El material vincula el muestreo con el modelo de riesgo de auditoría:

\begin{align}
    RA &= RI \times RC \times RD \\
    RD &= PA \times PD \\
    RA &= RI \times RC \times PA \times PD
\end{align}

donde $RA$ es el riesgo de auditoría, $RI$ el riesgo inherente, $RC$ el
riesgo de control, $RD$ el riesgo de detección, $PA$ el riesgo de no
detectar un error material con los procedimientos analíticos y $PD$ el
riesgo tolerable de aceptación incorrecta (que determina, en definitiva, el
nivel de riesgo con el que se entra a la Tabla~\ref{tab:riesgo}).

% ============================================================
\section{Descripción de los datos auditados}
% ============================================================

El listado \texttt{1393.csv} contiene 8.855 filas exportadas. Antes de
definir la población se eliminaron: (i) una fila vacía inicial y (ii) la
fila de totales del reporte (que no es un comprobante sino la suma de la
columna \emph{Total}). La suma de la columna \emph{Total} de la población
depurada se validó contra ese total de control, verificando el cuadre.

\begin{table}[htbp]
    \centering
    \caption{Estadísticos descriptivos de la variable \emph{Total} (Gs.)}
    \label{tab:descriptivas}
    \begin{tabular}{@{}lr@{}}
        \toprule
        Estadístico & Valor \\
        \midrule
        $N$ (población) & 8.853 \\
        Media & 1.526.889 \\
        Desviación estándar & 12.623.325 \\
        Mínimo & 0 \\
        Mediana & 593.600 \\
        Máximo & 1.166.521.726 \\
        Asimetría (skewness) & 88{,}85 \\
        \bottomrule
    \end{tabular}
\end{table}

La asimetría es muy alta porque un único comprobante (Fact. Crédito
001-001-0220680, del 06/02/2024, por \Gs{1.166.521.726}) concentra un
importe muy superior al resto (más de 38 desviaciones estándar por encima
de la media). Por ese motivo se decidió \textbf{examinar esa partida al
100\,\%} y aplicar el muestreo estadístico únicamente al resto de la
población ($N = \text{8.852}$ comprobantes), siguiendo el criterio de combinar
selección dirigida de partidas significativas con muestreo probabilístico
para el resto.

% ============================================================
\section{Selección del modelo de muestreo}
% ============================================================

La variable a auditar (\emph{Total}) es un monto, no un atributo
binario, por lo que corresponde aplicar el modelo de
\textbf{estimación de la medida por unidad} y no las fórmulas de Kish
(éstas se reservan para atributos, por ejemplo si el objetivo fuera medir
qué porcentaje de comprobantes carece de un dato obligatorio).

Los parámetros de riesgo se fijaron así: % TODO: justificar/ajustar según la materialidad real del trabajo

\begin{itemize}[nosep]
    \item Materialidad (error tolerable): 1\,\% del monto total de la
          población a muestrear.

    \item Nivel de riesgo aceptable: 5\,\% $\Rightarrow$ $CAI = 1{,}64$,
          $CRI = 1{,}96$ (Tabla~\ref{tab:riesgo}).
\end{itemize}

% ============================================================
\section{Cálculo del tamaño de muestra}
% ============================================================

Con la población depurada de la partida de alto valor:
$N = \text{8.852}$, $De = \Gs{2.451.313}$.

Error tolerable:

\[
    ET = 0{,}01 \times \Gs{12.351.024.537}
       = \Gs{123.510.000}
\]

Tolerancia planeada para el riesgo de muestreo:

\[
    TPRM = \frac{ET}{1 + \dfrac{CAI}{CRI}}
         = \frac{\text{123.510.000}}
                {1 + \dfrac{1{,}64}{1{,}96}}
         = \Gs{67.244.333{,}33}
\]

Tamaño de muestra:

\[
    n = \frac{N \times CRI \times De}{TPRM}
      = \frac{\text{8.852} \times 1{,}96 \times
              \text{2.451.313}}
             {\text{67.244.333{,}33}}
      = 632{,}47
      \;\Rightarrow\;
      n = 633
\]

% ============================================================
\section{Distribución de la muestra por estrato}
% ============================================================

Factor multiplicador: $f = n/N = 633/\text{8.852} = 7{,}15\%$.

\begin{table}[htbp]
    \centering
    \caption{Distribución de la muestra por tipo de documento}
    \label{tab:estratos}
    \begin{tabular}{@{}lrrr@{}}
        \toprule
        Tipo de documento & $N_h$ & $f$ & $n_h$ \\
        \midrule
        Factura de contado & 1.663 & 7{,}15\,\% & 119 \\
        Factura de crédito & 6.486 & 7{,}15\,\% & 464 \\
        Nota de crédito     &   703 & 7{,}15\,\% &  51 \\
        \midrule
        \textbf{Total}      & \textbf{8.852} & & \textbf{634} \\
        \bottomrule
    \end{tabular}
\end{table}

Dentro de cada estrato la selección se realizó al azar (muestreo aleatorio
simple estratificado), con semilla fija para que el resultado sea
reproducible.

% ============================================================
\section{Ficha técnica del muestreo}
% ============================================================

\begin{longtable}{@{}p{7cm}p{7cm}@{}}
    \caption{Ficha técnica del muestreo}
    \label{tab:ficha}\\
    \toprule
    \textbf{Parámetro} & \textbf{Valor} \\
    \midrule
    \endfirsthead

    \toprule
    \textbf{Parámetro} & \textbf{Valor} \\
    \midrule
    \endhead

    \midrule
    \multicolumn{2}{r}{\footnotesize\textit{Continúa en la siguiente página}} \\
    \endfoot

    \bottomrule
    \endlastfoot

    Población total ($N_0$) & 8.853 comprobantes \\
    Partidas de alto valor (examen 100\,\%) & 1 \\
    Población base de muestreo ($N$) & 8.852 comprobantes \\
    Variable de interés & Total del comprobante (Gs.) \\
    Modelo estadístico aplicado & Estimación de la medida por unidad \\
    Desviación estándar estimada ($De$) & \Gs{2.451.313} \\
    Nivel de riesgo aceptable & 5\,\% \\
    Coeficiente de aceptación incorrecta ($CAI$) & 1{,}64 \\
    Coeficiente de rechazo incorrecto ($CRI$) & 1{,}96 \\
    Error tolerable ($ET$) & \Gs{123.510.000} \\
    Tolerancia planeada para el riesgo de muestreo ($TPRM$) &
    \Gs{67.244.333{,}33} \\
    Tamaño de muestra calculado ($n$) & 633 \\
    Tipo de muestreo & Probabilístico, estratificado por tipo de documento \\
    Muestra final (alto valor + estadística) & 635 comprobantes \\
    Cobertura sobre la población total & 7{,}17\,\% \\
    Semilla de aleatorización & 42 \\

\end{longtable}

% ============================================================
\section{Conclusiones}
% ============================================================

% TODO: redactar las conclusiones propias del trabajo, por ejemplo:

\begin{itemize}
    \item El listado depurado quedó conformado por 8.853 comprobantes, cuya
    suma de importes coincide exactamente con el total de control del
    reporte original, lo que valida la integridad de la población.

    \item Dado que la variable auditada es un monto continuo, se aplicó el
    modelo de \emph{estimación de la medida por unidad}, obteniéndose un
    tamaño de muestra de 633 comprobantes sobre la población depurada del
    valor atípico.

    \item Sumando la partida de alto valor examinada al 100\,\%, la muestra
    final asciende a 635 comprobantes (7{,}17\,\% de la población total),
    distribuidos proporcionalmente entre los tres tipos de documento.
\end{itemize}

% ============================================================
% BIBLIOGRAFÍA
% ============================================================

\newpage

\begin{thebibliography}{9}

\bibitem{cofae}
COFAE -- Centro de Investigación y Desarrollo.
\emph{Muestreo en Auditoría}. Material de la cátedra CID Auditoría de
Estado. Docente adjunto: Pierre Ponte León.

\end{thebibliography}

% ============================================================
\end{document}